% Danh mục tài liệu tham khảo. Trích dẫn trong bài bằng \cite{khóa}; số thứ tự tự đánh.
\begin{thebibliography}{99}

\bibgroup{Văn bản pháp luật}
\bibitem{luat-attm-2015} Quốc hội (2015). \textit{Luật An toàn thông tin mạng}, số 86/2015/QH13. \href{https://hethongphapluat.com/luat-an-toan-thong-tin-mang-2015/dieu-7}{Điều 7}.
\bibitem{blhs-2015} Quốc hội (2015, sửa đổi 2017). \textit{Bộ luật Hình sự}, số 100/2015/QH13. \href{https://hethongphapluat.com/bo-luat-hinh-su-2015/dieu-288}{Điều 288}.
\bibitem{luat-anm-2018} Quốc hội (2018). \textit{Luật An ninh mạng}, số 24/2018/QH14. \href{https://hethongphapluat.com/luat-an-ninh-mang-2018/dieu-26}{Điều 26}.
\bibitem{nd53-2022} Chính phủ (2022). \textit{Nghị định 53/2022/NĐ-CP quy định chi tiết một số điều của Luật An ninh mạng}. Tóm tắt tại \href{https://lsvn.vn/tu-01102022-du-lieu-nguoi-dung-internet-tai-viet-nam-se-duoc-luu-tru1661052929-a122807.html}{LSVN}.
\bibitem{nd13-2023} Chính phủ (2023). \textit{Nghị định 13/2023/NĐ-CP về bảo vệ dữ liệu cá nhân}. \href{https://hethongphapluat.com/nghi-dinh-13-2023-nd-cp-ve-bao-ve-du-lieu-ca-nhan/khoan-3-dieu-2}{Điều 2 khoản 3}, \href{https://hethongphapluat.com/nghi-dinh-13-2023-nd-cp-ve-bao-ve-du-lieu-ca-nhan/dieu-3}{Điều 3}.
\bibitem{luat-bvdlcn-2025} Quốc hội (2025). \textit{Luật Bảo vệ dữ liệu cá nhân}, số 91/2025/QH15. \href{https://vanban.chinhphu.vn/?docid=214590&pageid=27160}{Hồ sơ chính thức}; \href{https://congbaocdn.chinhphu.vn/CongBaoCP/CongBao/2025/7/45574/57723-1-971-972.pdf}{bản Công báo, Điều 7, 8, 23}; mức phạt tại \href{https://lsvn.vn/mua-ban-du-lieu-ca-nhan-trai-phep-co-the-bi-phat-toi-10-lan-doanh-thu-a160473.html}{LSVN}.
\bibitem{nd356-2025} Chính phủ (2025). \textit{Nghị định 356/2025/NĐ-CP quy định chi tiết Luật Bảo vệ dữ liệu cá nhân}. \href{https://congbao.cdnchinhphu.vn/180507251028987904/2026/1/17/356signed-1768638052103952849513.pdf}{Bản Công báo}; tóm tắt tại \href{https://luatvietan.vn/diem-moi-cua-nghi-dinh-356-2025-so-voi-nghi-dinh-13-2023-ve-huong-dan-bao-ve-du-lieu-ca-nhan.html}{Luật Việt An}.
\bibitem{luat-qc-2012} Quốc hội (2012). \textit{Luật Quảng cáo}, số 16/2012/QH13. \href{https://hethongphapluat.com/luat-quang-cao-2012/dieu-8}{Điều 8}; \href{https://hrspring.vn/van-ban-luat/luat-quang-cao-sua-doi-2025-so-75-2025-qh15}{Luật sửa đổi số 75/2025/QH15}.
\bibitem{nd52-2013} Chính phủ (2013, sửa đổi 2021). \textit{Nghị định 52/2013/NĐ-CP về thương mại điện tử}, sửa đổi bởi Nghị định 85/2021/NĐ-CP. Phân tích tại \href{https://lsvn.vn/bao-dam-quyen-loi-cua-nguoi-tieu-dung-khi-mua-hang-tren-san-giao-dich-thuong-mai-dien-tu-1707211306-a140810.html}{LSVN}.
\bibitem{vbhn-tmdt-2020} Bộ Công Thương (2020). \textit{Văn bản hợp nhất quy định về thương mại điện tử}, \href{https://congbao.chinhphu.vn/tai-ve-van-ban-so-20-vbhn-bct-30960-30438?format=pdf}{Công báo 2020}, Điều 4.
\bibitem{nd52-2024} Chính phủ (2024). \textit{Nghị định 52/2024/NĐ-CP về thanh toán không dùng tiền mặt}. \href{https://vbpl.vn/nganhangnhanuoc/Pages/vbpq-toanvan.aspx?ItemID=167087}{Toàn văn}; \href{https://pbgdpl.cantho.gov.vn/mot-so-hanh-vi-bi-cam-trong-thanh-toan-khong-dung-tien-mat-tu-ngay-0172024}{Điều 8}.
\bibitem{luat-anm-2025} Quốc hội (2025). \textit{Luật An ninh mạng}, số 116/2025/QH15. \href{https://vanban.chinhphu.vn/?classid=1&docid=216499&pageid=27160&typegroupid=3}{Hồ sơ chính thức}; \href{https://apolatlegal.com/vi/laws/luat-an-ninh-mang-2025-so-116-2025-qh15/}{Apolat Legal}; \href{https://hethongphapluat.com/luat-an-ninh-mang-2025/dieu-44}{Điều 44}.
\bibitem{luat-tmdt-2025} Quốc hội (2025). \textit{Luật Thương mại điện tử}, số 122/2025/QH15. \href{https://luatvietnam.vn/tin-van-ban-moi/da-co-luat-thuong-mai-dien-tu-2025-so-122-2025-qh15-186-106402-article.html}{LuatVietnam}.
\bibitem{luat-ai-2025} Quốc hội (2025). \textit{Luật Trí tuệ nhân tạo}, số 134/2025/QH15, hiệu lực từ 1/3/2026. \href{https://thuvienphapluat.vn/van-ban/EN/Cong-nghe-thong-tin/Law-134-2025-QH15-Artificial-Intelligence/689735/tieng-anh.aspx}{Bản tiếng Anh}; tóm tắt tại \href{https://apolatlegal.com/vi/laws/luat-tri-tue-nhan-tao-2025-so-134-2025-qh15/}{Apolat Legal}.
\bibitem{nd142-2026} Chính phủ (2026). \textit{Nghị định 142/2026/NĐ-CP hướng dẫn Luật Trí tuệ nhân tạo}, hiệu lực từ 1/5/2026. Tóm tắt tại \href{https://luatvietan.vn/nghi-dinh-142-2026-nd-cp-huong-dan-luat-tri-tue-nhan-tao.html}{Luật Việt An}.

\bibgroup{Thông tin chính thức và báo chí về vụ taphoammo.net}
\bibitem{bca-2026} Bộ Công an (14/7/2026). \href{https://www.bocongan.gov.vn/bai-viet/triet-pha-nhom-doi-tuong-van-hanh-website-mua-ban-du-lieu-ca-nhan-tai-khoan-mang-xa-hoi-rac-1784022194}{Triệt phá nhóm đối tượng vận hành website mua bán dữ liệu cá nhân, tài khoản mạng xã hội}.
\bibitem{cand-2026} Báo Công an Nhân dân (9/7/2026). \href{https://cand.vn/mua-ban-du-lieu-duoi-vo-boc-cua-cong-ty-khoi-nghiep-3-doi-tuong-bi-khoi-to-post816257.html}{Mua bán dữ liệu dưới vỏ bọc của công ty khởi nghiệp, 3 đối tượng bị khởi tố}.
\bibitem{vnn-2026} VietNamNet (10/7/2026). \href{https://vietnamnet.vn/cong-an-tp-ha-noi-khoi-to-3-nguoi-dieu-hanh-san-mua-ban-du-lieu-ca-nhan-2534474.html}{Công an TP Hà Nội khởi tố 3 người điều hành sàn mua bán dữ liệu cá nhân}.
\bibitem{hnm-2026} Hà Nội Mới (10/7/2026). \href{https://hanoimoi.vn/khoi-to-cac-doi-tuong-mua-ban-du-lieu-ca-nhan-thu-loi-tren-400-ty-dong-1211456.html}{Khởi tố các đối tượng mua bán dữ liệu cá nhân thu lời trên 400 tỷ đồng}.
\bibitem{k14-2026} Kenh14 (11/7/2026, tổng hợp từ Đời sống \& Pháp luật). \href{https://kenh14.vn/hang-tram-canh-sat-ap-vao-kham-xet-tai-3-tinh-thanh-khoi-to-giam-doc-dao-quang-huy-sn-1992-215260710160532629.chn}{Hàng trăm cảnh sát ập vào khám xét tại 3 tỉnh thành, khởi tố Giám đốc Đào Quang Huy}.
\bibitem{bpl-2026} Báo Pháp luật Việt Nam (14/7/2026). \href{https://baophapluat.vn/triet-pha-website-taphoammo-net-ban-trai-phep-du-lieu-ca-nhan-thu-loi-tu-he-thong-1-35-trieu-tai-khoan.html}{Triệt phá website taphoammo.net bán trái phép dữ liệu cá nhân}.
\bibitem{dtud-2026} Tạp chí Điện tử và Ứng dụng (14/7/2026). \href{https://dientuungdung.vn/vu-tap-hoa-mmo-he-lo-mang-toi-mua-ban-du-lieu-ca-nhan-16575.html}{Vụ Tạp Hóa MMO hé lộ mảng tối mua bán dữ liệu cá nhân}.

\bibgroup{Hướng dẫn của cơ quan nhà nước}
\bibitem{cuc-tmdt-2026} Cục Thương mại điện tử (14/10/2026). \href{https://www.idea.gov.vn/default.aspx?category_id=6e3f3b28-3917-4816-83c5-a331259dd181&do=detail&id=184e0d3f-daab-42a5-a33e-782c78e579c0&page=news}{Thông báo về Luật Thương mại điện tử 122/2025/QH15, Nghị định 248/2026 và mốc xác thực danh tính}.
\bibitem{bca-huongdan-2026} Bộ Công an (2026). \href{https://bocongan.gov.vn/chinh-sach-phap-luat/bai-viet/canh-bao-nguy-co-vi-pham-phap-luat-ve-an-ninh-mang-doi-voi-doanh-nghiep-ho-kinh-doanh-trong-qua-trinh-chuyen-doi-so-1785806444}{Cảnh báo nguy cơ vi phạm pháp luật về an ninh mạng đối với doanh nghiệp, hộ kinh doanh trong quá trình chuyển đổi số}.

\bibgroup{Bối cảnh chung}
\bibitem{sggp-2021} Sài Gòn Giải Phóng (17/5/2021). \href{https://www.sggp.org.vn/triet-pha-duong-day-mua-ban-trai-phep-gan-1300gb-du-lieu-post594468.html}{Triệt phá đường dây mua bán trái phép gần 1.300GB dữ liệu}.
\bibitem{vneconomy-2024} VnEconomy (17/7/2024). \href{https://vneconomy.vn/phat-hien-hang-nghin-gb-du-lieu-ca-nhan-bi-thu-thap-mua-ban-trai-phep-trong-do-co-nhieu-du-lieu-nhay-cam.htm}{Phát hiện hàng nghìn GB dữ liệu cá nhân bị thu thập, mua bán trái phép}.
\bibitem{sggp-2026-thiethai} Sài Gòn Giải Phóng (7/1/2026). \href{https://www.sggp.org.vn/nam-2025-thiet-hai-do-lua-dao-truc-tuyen-o-viet-nam-uoc-tinh-tren-6000-ty-dong-post832336.html}{Năm 2025, thiệt hại do lừa đảo trực tuyến ở Việt Nam ước tính trên 6.000 tỷ đồng}.
\bibitem{sggp-2026-50trieu} Sài Gòn Giải Phóng (7/7/2026). \href{https://www.sggp.org.vn/bat-2-doi-tuong-mua-ban-trai-phep-hon-50-trieu-du-lieu-ca-nhan-post861086.html}{Bắt 2 đối tượng mua bán trái phép hơn 50 triệu dữ liệu cá nhân}.
\end{thebibliography}
